% Add-ons: adapters that read recording formats Medorama and the viewers do not read themselves.
% What they are, installing them (medorama addons install), how the programs use them, the
% adapter protocol in full, writing one, security, and the add-ons known (the STP add-on).
% Sources: docs/adapters.md; medorama/src/addons.rs, main.rs (addons); libs/engine/src/adapters.rs
% (PROTOCOL = 1, describe, default_folder).
\chapter{Add-ons}
\label{chap:addons}

This chapter explains adapters: separate programs that let Medorama and the VitalStream viewers
read recording formats they do not read themselves. It is written for administrators who install
an add-on, and for developers who write an adapter for a new format.

Medorama and the VitalStream viewers read WFDB records and EDF, EDF+, BDF and BDF+ files
themselves (\cref{chap:recordings}). Every other recording format is read by an
\textbf{adapter}: a program, installed beside them, that turns a recording into the index the
viewers draw from (\cref{chap:formats-reference}), and returns its samples when a view zooms in.
Formats that are licensed, or that only some institutions may read, are distributed this way, as
\textbf{add-ons}, outside the public repository. When to write an adapter rather than a built-in
reader is discussed in \cref{sec:readers-recordings}. This chapter describes what an adapter is,
how an administrator installs one, how the programs use it, the protocol that an adapter speaks,
and the add-ons known.

\section{What an adapter is}
\label{sec:addons-what}

\begin{xltabular}{\linewidth}{@{}P{3.0cm}L@{}}
\caption{An adapter.}
\label{tab:addons-adapter}\\
\toprule
\tabhead Aspect & \tabhead An adapter \\
\midrule
\endfirsthead
\toprule
\tabhead Aspect & \tabhead An adapter \\
\midrule
\endhead
\bottomrule
\endfoot
A program & An executable named \file{<name>-adapter} (\file{stp-adapter}; on Windows,
  \file{stp-adapter.exe}). \\
Where & An \textbf{adapter folder}: \file{\textasciitilde/.local/share/medorama/adapters} by
  default (on Windows, \code{\%LOCALAPPDATA\%\textbackslash medorama\textbackslash adapters}).
  The setting \code{adapters\_dir} of Medorama (\file{medorama.toml}) and of the viewers
  (\file{vitalstream.toml}, or \opt{-{}-adapters-dir}) names another. \\
What it reads & The formats it lists when asked (\code{describe}), each with a recording kind, a
  label and file extensions. \\
What it writes & Only the index files, into the folder it is given; never anything next to a
  recording. \\
How it is called & By Medorama and by each viewer, as a child process, with the commands of
  \cref{sec:addons-protocol}. Nothing runs as a server or listens on a port. \\
\end{xltabular}

Medorama and the viewers name no add-on's format. Until its adapter is installed, a file of that
format is treated like any other file they do not read: the file browsers do not list it, and
opening it is refused (\code{not a recording this program reads}). The built-in readers take
precedence over an adapter that lists the same extensions (\file{.hea}, \file{.edf},
\file{.bdf}).

\section{Installing an add-on}
\label{sec:addons-installing}

An add-on arrives as an archive for one system,
\file{<name>-adapter-<version>-<os>-<arch>.tar.gz} (Linux, macOS) or \file{.zip} (Windows),
holding the program, its license and \file{INSTALL.txt}. An administrator installs it on the
server, never through the page, because an adapter is a program that runs with the catalog's
rights.

\paragraph{1. Unpack it, and allow the system to run it.} Programs downloaded from the internet
are blocked until you allow them. The following commands unpack the archive on Linux and macOS,
and remove macOS's quarantine mark (otherwise Gatekeeper refuses to run the program):
\begin{lstlisting}[style=shell]
tar -xzf stp-adapter-0.1.0-linux-x86_64.tar.gz                  # Linux and macOS
xattr -d com.apple.quarantine stp-adapter-0.1.0-*/stp-adapter   # macOS only (Gatekeeper)
\end{lstlisting}
On Windows, the following commands unpack the archive and unblock the program:
\begin{lstlisting}[language=PowerShell]
Expand-Archive stp-adapter-0.1.0-windows-amd64.zip .
Unblock-File stp-adapter-0.1.0-windows-amd64\stp-adapter.exe
\end{lstlisting}

\paragraph{2. Install it, as the account that runs the catalog} (or with \opt{-{}-home} naming
the catalog), so that it goes into that catalog's adapter folder (\code{adapters\_dir}). The
following commands install it and list the installed adapters:
\begin{lstlisting}[style=shell]
medorama addons install stp-adapter-0.1.0-linux-x86_64/stp-adapter
medorama addons list            # name, version, implementation, formats, program
\end{lstlisting}
\cmd{addons install} checks the program first: it must be named \file{<name>-adapter}
(\file{.exe} on Windows), answer \code{describe}, and speak this protocol's version. It then
copies the program into the adapter folder (creating the folder if needed; on Unix with mode
0755), replacing an installed adapter of the same name, and asks the copy to describe itself
again. Installing and removing are written to the audit log (\code{addon\_install},
\code{addon\_remove}). \cmd{medorama addons remove stp} removes an adapter by name.

\paragraph{3. Restart Medorama} (\cmd{systemctl -{}-user restart medorama}, or stop and start
\cmd{medorama serve}). The programs list their adapters when they start. Stopping Medorama also
stops the VitalStream viewer it runs, so the next viewer sees the add-on too. \ui{Admin}
$\rightarrow$ \ui{Add-ons} then lists it.

\paragraph{4. Import again the tables that named such files before.} An import skips a recording
that no installed reader reads, and that no viewer opens either (\code{line N: FILE: not a
recording Medorama reads. A format read by an add-on needs it installed ...}), and keeps nothing
of it. Once the add-on is installed, remove that earlier import (\ui{Import} $\rightarrow$
\ui{Sources}, or \cmd{medorama imports list} and then \cmd{medorama imports remove ID}) and
import the table again. The files then come in as recordings, linked to their patients and
indexed by the workers. Importing the table again without removing the first import makes a
second copy of the database.

\paragraph{The viewers on their own.} Without Medorama, copy the program into the viewers'
adapter folder (by default the same, \file{\textasciitilde/.local/share/medorama/adapters}, or
their \code{adapters\_dir} setting), make it executable on Linux and macOS (\cmd{chmod 755}),
unblock it as in step 1, and restart the viewer.

\begin{xltabular}{\linewidth}{@{}P{5.6cm}L@{}}
\caption{Problems when installing an add-on.}
\label{tab:addons-problems}\\
\toprule
\tabhead Problem & \tabhead Cause and remedy \\
\midrule
\endfirsthead
\toprule
\tabhead Problem & \tabhead Cause and remedy \\
\midrule
\endhead
\bottomrule
\endfoot
\cmd{install}: \code{an adapter's program is named <name>-adapter} & The file was renamed.
  Install the program as it came in the archive. \\
\cmd{install}: \code{describe} fails, ``cannot be opened'', ``permission denied'' & The system
  blocks the program. Repeat step 1 (\cmd{xattr} on macOS, \cmd{Unblock-File} on Windows), or
  run \cmd{chmod 755}. \\
\cmd{install}: \code{speaks protocol N; this program speaks 1} & The add-on was built for another
  Medorama version. Ask for the release that matches yours. \\
\cmd{addons list} shows it, but its files are still refused & The server was not restarted
  (step 3), or it runs as another account whose \code{adapters\_dir} is elsewhere. Run
  \cmd{medorama addons list} as that account. \\
Recordings named by an earlier import are missing & They were skipped before the add-on was
  installed. Remove that import and import the table again (step 4). \\
\end{xltabular}

\section{How the programs use adapters}
\label{sec:addons-use}

\begin{enumerate}
\item \textbf{At start}, each program lists its adapter folder and asks every
  \file{*-adapter} program to \code{describe} itself, once. When two adapters have the same
  name, the first in name order is used.
\item \textbf{Browsing.} A file whose extension an adapter lists is a recording of that adapter's
  kind.
\item \textbf{Indexing.} The program creates a temporary index folder, runs
  \code{index FILE FOLDER} and shows the progress lines. It then adds the fields it owns
  (\code{path}, \code{name}, \code{key}, \code{index\_version}, \code{engine},
  \code{indexed\_s}, and \code{kind}, \code{format}, \code{adapter}) to the \file{meta.json}
  that the adapter wrote, and renames the folder into place.
\item \textbf{Viewing.} Envelopes, markers and vital signs come from the index. Raw samples, when
  a view zooms in close enough, are read from the recording itself through the adapter's
  \code{serve} session. There is one session per adapter, started at the first request and kept
  open, and started again once if it has ended. The sessions of different adapters answer side
  by side; one adapter's requests are answered in turn.
\item \textbf{Identifiers} (Medorama only). The built-in action \code{identify} asks the adapter
  for the identifiers the file states (\code{identify FILE}), and keeps them in the catalog's
  database only (\cref{chap:administration}).
\end{enumerate}
Indexes written through an adapter are shared like the others: the three viewers and Medorama
read and write the same index folders, and the three implementations of the stp-adapter write
the same bytes.

\section{The protocol}
\label{sec:addons-protocol}

This is version 1 of the protocol. Every command writes its answer on the standard output as
JSON (UTF-8, one object per line), and its errors on the standard error, with a nonzero exit
status. Paths are given as
arguments, as the operating system names them.

\apientry{describe}{What the adapter reads}
\begin{lstlisting}[style=signature]
<name>-adapter describe
\end{lstlisting}
Prints one object: the adapter's name, its version, the protocol version it speaks, an
informative \code{implementation}, and its formats, each a recording kind, a label and the file
extensions (without the dot, in lowercase). A program refuses an adapter whose \code{protocol} is
not 1.
\begin{lstlisting}[language=JSON]
{"name": "stp", "version": "0.1.0", "protocol": 1, "implementation": "rs",
 "formats": [{"kind": "stp", "label": "BedMasterEx .Stp", "extensions": ["stp"]}]}
\end{lstlisting}

\apientry{detect}{Whether files are of its formats}
\begin{lstlisting}[style=signature]
<name>-adapter detect PATH...
\end{lstlisting}
Prints one line per path: \code{\{"path": ..., "kind": "stp"\}}, with \code{kind} \code{null}
when the file is not of its formats.

\apientry{index}{Index a recording}
\begin{lstlisting}[style=signature]
<name>-adapter index PATH DIR
\end{lstlisting}
Writes the index files of the recording at \code{PATH} into the existing, empty folder
\code{DIR}. It prints \code{\{"progress": f\}} lines (\code{f} from 0 to 1, a few per second at
most), and then \code{\{"done": true\}}.

The index follows the index format (\cref{chap:formats-reference}) in everything but the fields
the calling program adds (\cref{sec:addons-use}, step 3): a \file{meta.json} with
\code{channels}, \code{pieces}, \code{vitals}, \code{alarms}, \code{gaps}, \code{untimed},
\code{marker\_names}, \code{skipped}, \code{notes}, \code{time\_resolution\_s} and
\code{reader}, and every array file. An adapter may add keys of its own (\code{vendor}, and
\code{description}, a sentence shown with the recording), and keep in \code{reader} whatever it
needs to read samples later. The programs never interpret \code{reader}. An adapter that puts
data the file stored out of order back in order lists those blocks in \code{late\_blocks}: the
viewers do not shade them as breaks, and mark them only with their \code{qc\_verbose} setting.

\apientry{identify}{The identifiers a file states}
\begin{lstlisting}[style=signature]
<name>-adapter identify PATH...
\end{lstlisting}
Prints one line per path, \code{\{"path": ..., "identifiers": \{...\} or null, "others": [\{...\},
...]\}}, or \code{\{"path": ..., "error": "..."\}}. A set of identifiers has the fields
\code{full\_name}, \code{last\_name}, \code{first\_name}, \code{patient\_id},
\code{encounter\_id}, \code{birth\_date} (\code{YYYY-MM-DD}) and \code{sex}, each a string,
empty when the file does not state it. \code{identifiers} is the first set the file states, and
\code{others} lists the different sets found after it (at most 10). Only Medorama asks for
identifiers. The example below is invented:
\begin{lstlisting}[language=JSON]
{"path": "/data/demo/rec_0001.stp", "identifiers": {"full_name": "TEST, PATIENT",
 "last_name": "TEST", "first_name": "PATIENT", "patient_id": "SYN0001", "encounter_id": "",
 "birth_date": "1950-01-01", "sex": "F"}, "others": []}
\end{lstlisting}

\apientry{serve}{A session for samples}
\begin{lstlisting}[style=signature]
<name>-adapter serve
\end{lstlisting}
A session on the standard input and output, until the standard input closes. Each request is one
line of JSON:
\begin{lstlisting}[language=JSON]
{"op": "read", "index": "/path/to/the/index/folder", "piece": 3, "a": 1000, "b": 3500}
\end{lstlisting}
It asks for the physical samples \code{a} to \code{b - 1} of piece \code{piece} of the recording
indexed in \code{index}, whose \file{meta.json} now holds the recording's \code{path}. The answer
is one line, \code{\{"n": k\}}, followed at once by \code{k} little-endian 64-bit floats (NaN for
the samples the recording marks missing), or one line \code{\{"error": "..."\}} with nothing
after it. An adapter answers requests in order, may keep recordings open between them, and ends
when its standard input closes.

\apientry{-{}-version}{The adapter's version}
\begin{lstlisting}[style=signature]
<name>-adapter --version
\end{lstlisting}
Prints \code{<program> <version>} on one line.

\section{Limits and environment}
\label{sec:addons-limits}

The calling program holds every adapter to the limits of \cref{tab:addons-limits}, so that a
faulty or stuck adapter makes one recording unreadable and never holds up the program. The Rust
engine (\file{libs/engine/src/adapters.rs}) defines them; the C++ and Python viewers apply the
same.

\begin{xltabular}{\linewidth}{@{}P{1.6cm}P{7.0cm}L@{}}
\caption{The limits on an adapter's commands.}
\label{tab:addons-limits}\\
\toprule
\tabhead Command & \tabhead Limit & \tabhead When it is exceeded \\
\midrule
\endfirsthead
\toprule
\tabhead Command & \tabhead Limit & \tabhead When it is exceeded \\
\midrule
\endhead
\bottomrule
\endfoot
\code{describe} & 30\,s for the whole command; at most 16\,MiB of output & The adapter is
  stopped (killed) and not listed \\
\code{identify} & 60\,s per path of the command; at most 16\,MiB of output. The paths are passed
  in commands whose arguments total at most 24\,000 bytes (one path per command at least) & The
  adapter is stopped; the identifiers are not read \\
\code{index} & 300\,s without a line of output (a progress line resets the time), and 300\,s for
  the adapter to exit once its output ends; lines of at most 1\,MiB. A cancelled index stops the
  adapter within a tenth of a second & The adapter is stopped; the index fails, and nothing is
  left in the index folder \\
\code{serve} & 60\,s without a byte of an answer; an answer line of at most 1\,MiB. The program
  asks for at most $2^{20}$ samples in one request and splits longer reads. The answer's
  \code{n} must be a whole number from 0 to \code{b - a}; fewer samples than asked mean that the
  piece ends there & The session is stopped (below) \\
\end{xltabular}

\paragraph{A session out of step is stopped.} After an answer that breaks the protocol (not a
JSON object, an \code{n} outside \code{0 \dots{} b - a}, bytes beyond the samples announced, the
session ending inside an answer, or no answer in time), the program stops the session and starts
a new one at the next request; the request fails. An answer \code{\{"error": "..."\}} keeps the
session. A session that has ended before a request is started again once for that request.

\paragraph{Standard error} is read while the adapter runs, so an adapter that writes much there
never blocks; the program keeps its last 64\,KiB and never prints them in its log. When
\code{index} fails, the error shown with the recording ends with the last non-blank line of the
adapter's standard error (at most 200 characters, control characters removed); when
\code{describe} fails, the program's log shows the same line. What \code{identify} and
\code{serve} write there is not shown at all, since it could carry identifiers.

\paragraph{The environment} of an adapter is cleared, except for \code{PATH}, \code{HOME},
\code{USERPROFILE}, \code{TMP}, \code{TEMP}, \code{TMPDIR}, \code{LANG}, the \code{LC\_*}
variables, and on Windows \code{SYSTEMROOT}, \code{SYSTEMDRIVE}, \code{WINDIR}, \code{COMSPEC},
\code{PATHEXT}, \code{LOCALAPPDATA} and \code{APPDATA} (names compared ignoring case on
Windows). Medorama keeps secrets in its environment (the database password, its key); an adapter
never sees them. An adapter must not depend on any other variable: what it needs to know comes
in its arguments and requests.

\section{Writing an adapter}
\label{sec:addons-writing}

An adapter can be written in any language. The engines' building blocks make it short in the
three languages that the repository uses:
\begin{itemize}
\item \textbf{Rust.} The crates \code{recording-engine} (\code{common}, with channels, pieces,
  envelopes and vital-sign series; and \code{index::write\_arrays}) and \code{signal-core}
  (values and Python-compatible JSON). The test adapter of \file{libs/engine/tests/adapters.rs}
  is an example.
\item \textbf{C++.} The C++ project's libraries \code{vs\_engine} and \code{vs\_core}
  (\file{engine/common.hpp}, and \file{engine/index.hpp} with \code{write\_arrays}). Add
  \file{viewers/vitalstream/cpp} to a CMake project with \code{VS\_BUILD\_VIEWER=OFF}
  (\cref{chap:settings-reference}).
\item \textbf{Python.} \code{vitalstream\_py.common} and
  \code{vitalstream\_py.index.write\_arrays}. \file{libs/engine/tests/fake\_adapter.py}, the mock
  adapter that every engine's tests use, is an example.
\end{itemize}
An adapter must meet these rules:
\begin{itemize}
\item It writes only into the folder that \code{index} is given, and nothing at all for the other
  commands.
\item It never prints identifiers except in answer to \code{identify} (on standard output; never
  on standard error), and never writes them into the index.
\item It answers within the limits (\cref{sec:addons-limits}).
\item Times are Unix seconds (UTC), as in the index format.
\item The same recording gives the same index bytes every time: the index holds no time stamps of
  the indexing and no absolute paths of temporary files.
\item It is tested with the viewers. \cmd{viewers/vitalstream/tests/parity/index\_parity.py
  -{}-adapters-dir DIR FILE...} indexes a recording with every built viewer through the adapter
  and checks that the results are byte-identical.
\end{itemize}

\section{Security}
\label{sec:addons-security}

\begin{itemize}
\item \textbf{An adapter is code that runs with the rights of the program that calls it} (the
  catalog's service account, or a viewer's user). Install only adapters from their publisher, and
  keep the adapter folder writable by its owner only (the default folder is in that user's home).
\item \textbf{It reads recordings, which may hold identifiers.} It must not send anything over the
  network, and it writes only into the index folder it is given. The stp-adapter opens no
  network connection and writes no other file.
\item \textbf{The identifiers it reads} (\code{identify}) go only to Medorama's database, as other
  identifiers do (\cref{chap:users-privacy}). The index holds none.
\item \textbf{A failing adapter} (a crash, or a malformed answer) makes the recording unreadable,
  never the program. The indexing fails with the adapter's message, and a session that ended is
  started again once. An adapter that stops answering is stopped after a time limit
  (\cref{sec:addons-limits}).
\item \textbf{It gets none of the program's secrets:} its environment is cleared except for a
  short allow-list, and its standard error never reaches the program's log.
\end{itemize}

\section{The add-ons known}
\label{sec:addons-known}

\paragraph{The STP add-on} (\file{stp-adapter}) reads the \file{.Stp} files of BedMasterEx,
which records the waveforms, vital signs and alarms of GE bedside monitors (Unity, Carescape) and
Philips bedside monitors (PIICiX, Viridia). It indexes them for the viewers, reads their samples
on demand through its \code{serve} session, and answers \code{identify} with the identifiers the
file states. It opens no network connection and writes no file other than the index. Its three
implementations (Rust, C++ and Python) write the same index bytes for the same file. It describes its
format as \code{BedMasterEx .Stp}, of kind \code{stp}, with the extension \code{stp}.

\begin{note}[License]
The STP add-on is private. It is licensed separately and released, as a prebuilt program for
their systems, to institutions under agreement. It is not part of the public repository, and
Medorama names none of its formats: until it is installed, \file{.Stp} files are files that
Medorama and the viewers do not read.
\end{note}
